\section{Heterogeneous mixed-radix tables}

\subsection{Why one natural number is enough}

At a universal node a certificate may need one subcertificate for every
\(i<N\).  The number of entries depends on the input, and the bound on the
\(i\)-th subcertificate may depend on earlier quantified values.  A fixed tuple
type is therefore inadequate, while an abstract list would require a separate
list theory and coding layer.

Natural numbers already provide the needed finite container.  Let
\(B_0,\ldots,B_{N-1}\) be positive radices and let
\(0\le c_i<B_i\).  Define prefix products and the code by
\[
  P_0=1,\qquad P_{i+1}=P_iB_i,
  \qquad
  \pack_N(\vec c)=\sum_{i<N}c_iP_i.
\]
The total code range is \(0\le\pack_N(\vec c)<P_N\).  The \(i\)-th digit is
recovered by
\[
  \digit_i(x)=\bigl(\lfloor x/P_i\rfloor\bmod B_i\bigr).
\]
All operations are primitive recursive when \(N\), the digit function, and the
radix function are primitive recursive.  In the repository they are assembled
from bounded sums, bounded products, quotient, and remainder.

\subsection{The interval lemma}

The basic estimate is elementary but central.

\begin{lemma}[Append interval]
If \(0\le x<P\), \(0\le d<B\), and \(P,B>0\), then
\[
  0\le x+dP<PB.
\]
\end{lemma}

\begin{proof}
The lower bound is immediate.  Since \(d<B\), natural-number discreteness gives
\(d+1\le B\).  Therefore
\[
  x+dP < P+dP=(d+1)P\le BP.
\]
Every step is an order or semiring identity already available in ordinary iPRA.
The proof uses no subtraction by an unknown sign and no classical case split.
\end{proof}

\begin{theorem}[Packed-code bound]
Assume \(B_i>0\) and \(c_i<B_i\) for every \(i<N\).  Then
\[
  \pack_N(\vec c)<P_N.
\]
\end{theorem}

\begin{proof}
Use induction on \(N\).  For \(N=0\), the empty sum is zero and the empty
product is one.  In the successor step,
\[
 \pack_{N+1}=\pack_N+c_NP_N,
 \qquad P_{N+1}=P_NB_N.
\]
The induction hypothesis gives \(\pack_N<P_N\), and the append interval lemma
with \(x=\pack_N\), \(P=P_N\), \(d=c_N\), and \(B=B_N\) gives the result.  The
induction predicate is an inequality between primitive-recursive terms and is
quantifier-free.
\end{proof}

\subsection{Digit recovery}

Fix \(i<N\).  Split the packed sum into lower, current, and upper pieces:
\[
  \pack_N=L_i+c_iP_i+P_{i+1}U_i,
\]
where
\[
  L_i=\sum_{j<i}c_jP_j,
  \qquad
  U_i=\sum_{i<j<N}c_j\prod_{i<k<j}B_k.
\]
By the packed-code bound applied to the prefix, \(L_i<P_i\).  Since
\(P_{i+1}=P_iB_i\), we may rewrite
\[
 \pack_N=P_i(c_i+B_iU_i)+L_i.
\]
Division uniqueness for the positive divisor \(P_i\) therefore gives
\[
  \left\lfloor\frac{\pack_N}{P_i}\right\rfloor=c_i+B_iU_i,
  \qquad
  \pack_N\bmod P_i=L_i.
\]
A second use of division uniqueness, now for divisor \(B_i\), gives remainder
\(c_i\).  Thus \(\digit_i(\pack_N)=c_i\).

\begin{theorem}[Mixed-radix decoding]
Under the hypotheses of the packed-code bound, every \(i<N\) satisfies
\[
  \digit_i(\pack_N(\vec c))=c_i.
\]
\end{theorem}

\begin{proof}
The preceding calculation is the mathematical proof.  For a direct iPRA
formalization it is convenient to avoid defining \(U_i\) first.  Prove by
induction on \(N\) two lemmas simultaneously.

\begin{enumerate}
\item The newest digit of \(x+dP_N\) is \(d\) whenever \(x<P_N\) and
\(d<B_N\).
\item Appending \(dP_N\) leaves every earlier digit unchanged.
\end{enumerate}

For the newest digit, division uniqueness by \(P_N\) identifies quotient
\(d\) and remainder \(x\), and division by \(B_N\) identifies the remainder of
\(d\) as \(d\).  For an earlier index \(i<N\), \(P_N\) is divisible by
\(P_iB_i\); adding a multiple of \(P_iB_i\) changes the quotient by a multiple
of \(B_i\) and hence does not change its remainder modulo \(B_i\).  The
repository's finite-product recurrence supplies the factorization of prefix
products.  Applying the induction hypothesis to the old prefix completes the
proof.
\end{proof}

\subsection{Uniqueness and extensionality}

\begin{corollary}[Code extensionality]
If \(c_i,d_i<B_i\) for all \(i<N\), then
\[
  \pack_N(\vec c)=\pack_N(\vec d)
  \quad\Longleftrightarrow\quad
  \forall i<N\;c_i=d_i.
\]
\end{corollary}

\begin{proof}
The reverse implication follows from bounded-sum congruence.  For the forward
implication, apply \(\digit_i\) to both sides and use the decoding theorem.  In
the object proof, ``for all \(i<N\)'' is represented by a Boolean fold or by a
universally quantified implication whose matrix is quantifier-free; the
compiler chooses the form appropriate to the consumer.
\end{proof}

\subsection{Disjoint sums for one selected child}

Tables encode a certificate for every child.  A node controlled by the winning
side needs only one child.  Let child code space \(i\) have size \(B_i\), which
may be zero when no certificate of the requested sign can exist.  Define prefix
sums
\[
  S_0=0,
  \qquad S_{i+1}=S_i+B_i.
\]
The disjoint union has bound \(S_N\).  A pair \((i,c)
equation\) with
\(i<N\) and \(c<B_i\) is encoded as
\[
  S_i+c.
\]
Conversely, for every \(x<S_N\), the least \(i<N\) satisfying
\(x<S_{i+1}\) exists; minimality gives \(S_i\le x\), and
\(c=x\mathbin{\dot{-}}S_i<B_i\).  The owner index is computed by bounded
least search.  This is a particularly clean use of the existing
\code{Core.Search}: it supplies both foundness and minimality.

\begin{lemma}[Disjoint decoding]
For \(i<N\) and \(c<B_i\), the least owner of \(S_i+c\) is \(i\), and the local
code is \(c\).
\end{lemma}

\begin{proof}
For every \(j<i\), monotonicity of prefix sums gives
\(S_{j+1}\le S_i\le S_i+c\), so no earlier interval contains the code.  At
\(i\), \(c<B_i\) gives \(S_i+c<S_i+B_i=S_{i+1}\).  Least-search minimality and
truncated-subtraction cancellation yield the local code equation.
\end{proof}

\subsection{Zero radices and empty tables}

A tight table bound \(\prod_{i<N}B_i\) can be zero.  This is correct: if one
required child has no certificate of the demanded sign, no table certificate
exists.  The checker first verifies \(c<\prod B_i\).  Under that assumption the
product is positive, hence every factor is positive, so quotient/remainder
lemmas may safely be used.  If the table is empty, the product is one and the
unique code is zero.

For an initial Lean implementation one may instead use padded radices
\(B_i+1\), reject the sentinel digit \(B_i\), and later optimize to the tight
bound.  STRIVE's reference implementation uses the tight representation, and
the formalization plan states the positivity obligations explicitly.

\subsection{Primitive-recursive realization}

Let \(b(i,p)\) and \(c(i,p)\) be concrete primitive-recursive terms.  Define
\[
 \pref(i,p)=\prod_{j<i}b(j,p),
 \qquad
 \pack(N,p)=\sum_{i<N}c(i,p)\pref(i,p).
\]
The current repository already supplies parameterized finite products and
sums.  Hence \(\pref\) and \(\pack\) are concrete \(PR\) symbols with zero and
successor equations.  Digit access is the term
\[
 \operatorname{rem}\bigl(
   \operatorname{quo}(x,\pref(i,p)),\;b(i,p)
 \bigr).
\]
The Lean blueprint \code{STRIVE/MixedRadix.lean} is organized precisely around
these existing equations.  No abstract sequence type, quotient type, or
universal decoding relation is required.
